\documentclass[aspectratio=169]{beamer}
\usetheme{Madrid}
\usecolortheme{whale}
\usepackage[utf8]{inputenc}
\usepackage[spanish]{babel}
\usepackage{amsmath,amsfonts,amssymb}
\usepackage{graphicx}
\usepackage{booktabs}
\usepackage{xcolor}

\title{Métodos de Análisis de Datos 1}
\subtitle{Clase 9: Análisis de Distribuciones}
\author{Depto.\ de Matemática -- UNS}
\institute{Universidad Nacional del Sur}
\date{1er Cuatrimestre 2026}

\AtBeginSection[]{
  \begin{frame}{Contenidos}
    \tableofcontents[currentsection]
  \end{frame}
}

\begin{document}

\begin{frame}
  \titlepage
\end{frame}

\begin{frame}{Contenidos}
  \tableofcontents
\end{frame}

%=============================================================
\section{Histogramas}
%=============================================================

\begin{frame}{Histograma: definición y construcción}
  \begin{block}{¿Qué es?}
    Representación gráfica de la distribución de frecuencias de una variable continua, mediante rectángulos adyacentes.
  \end{block}
  \bigskip
  \textbf{Construcción:}
  \begin{enumerate}
    \item Determinar el número de clases: regla de Sturges $k = 1 + 3.322\log_{10}(n)$.
    \item Calcular amplitud $h = (x_{\max}-x_{\min})/k$.
    \item Dibujar rectángulos donde el \textbf{área} es proporcional a la frecuencia relativa.
  \end{enumerate}
  \bigskip
  \begin{alertblock}{Área vs.\ altura}
    El eje vertical debe ser \textbf{densidad} ($f_i/h$) cuando las clases tienen diferente amplitud, no frecuencia.
  \end{alertblock}
\end{frame}

\begin{frame}{Formas típicas de histogramas}
  \begin{itemize}
    \item \textbf{Simétrico / campana}: datos centrados, colas equilibradas. Sugiere distribución Normal.
    \item \textbf{Sesgado a la derecha}: cola larga hacia valores altos. Ej.: ingresos, tiempos de falla.
    \item \textbf{Sesgado a la izquierda}: cola larga hacia valores bajos. Ej.: notas en examen fácil.
    \item \textbf{Bimodal}: dos picos — posible mezcla de subpoblaciones.
    \item \textbf{Uniforme}: todas las clases con frecuencia similar.
  \end{itemize}
  \bigskip
  \begin{exampleblock}{Interpretación}
    La forma del histograma orienta sobre qué distribución teórica podría ajustarse a los datos.
  \end{exampleblock}
\end{frame}

%=============================================================
\section{Boxplot}
%=============================================================

\begin{frame}{Diagrama de caja y bigotes (Boxplot)}
  \begin{block}{Componentes del boxplot}
    \begin{itemize}
      \item \textbf{Caja}: desde $Q_1$ hasta $Q_3$ (contiene el 50\% central).
      \item \textbf{Línea central}: mediana $Q_2$.
      \item \textbf{Bigote inferior}: hasta $Q_1 - 1.5 \cdot \text{RIC}$ (o el mínimo si no hay outliers).
      \item \textbf{Bigote superior}: hasta $Q_3 + 1.5 \cdot \text{RIC}$ (o el máximo si no hay outliers).
      \item \textbf{Puntos aislados}: observaciones fuera de los bigotes $\Rightarrow$ posibles outliers.
    \end{itemize}
  \end{block}
  \bigskip
  \begin{alertblock}{Ventajas}
    Muestra simultáneamente: posición central, dispersión, asimetría y outliers.
  \end{alertblock}
\end{frame}

\begin{frame}{Lectura e interpretación del boxplot}
  \textbf{Asimetría:}
  \begin{itemize}
    \item Si la mediana está centrada en la caja: distribución aproximadamente simétrica.
    \item Si la mediana está cercana a $Q_1$: sesgo positivo (cola derecha).
    \item Si la mediana está cercana a $Q_3$: sesgo negativo (cola izquierda).
  \end{itemize}
  \bigskip
  \textbf{Comparación de grupos:}
  \begin{itemize}
    \item Los boxplots son especialmente útiles para comparar la distribución de una variable cuantitativa entre varios grupos (variable cualitativa en el eje $x$).
  \end{itemize}
\end{frame}

%=============================================================
\section{Q-Q Plot}
%=============================================================

\begin{frame}{Gráfico cuantil-cuantil (Q-Q plot)}
  \begin{block}{Definición}
    Compara los cuantiles empíricos de los datos contra los cuantiles teóricos de una distribución de referencia (generalmente la Normal).
  \end{block}
  \bigskip
  \textbf{Construcción:}
  \begin{enumerate}
    \item Ordenar los datos: $x_{(1)} \le x_{(2)} \le \cdots \le x_{(n)}$.
    \item Calcular probabilidades empíricas: $p_i = (i - 0.5)/n$ (o variantes).
    \item Obtener cuantiles teóricos: $z_i = \Phi^{-1}(p_i)$.
    \item Graficar $(z_i, x_{(i)})$.
  \end{enumerate}
  \bigskip
  \begin{alertblock}{Interpretación}
    Si los puntos siguen la recta diagonal $\Rightarrow$ los datos son compatibles con la distribución teórica.
  \end{alertblock}
\end{frame}

\begin{frame}{Patrones en el Q-Q plot Normal}
  \begin{itemize}
    \item \textbf{Línea recta}: normalidad aceptable.
    \item \textbf{Curvatura hacia arriba}: sesgo positivo (cola derecha más pesada).
    \item \textbf{Curvatura hacia abajo}: sesgo negativo.
    \item \textbf{Forma de ``S''}: curtosis distinta a la Normal (colas más livianas o más pesadas).
    \item \textbf{Puntos alejados en los extremos}: outliers o colas pesadas.
  \end{itemize}
  \bigskip
  \begin{exampleblock}{Complemento}
    El Q-Q plot es una herramienta visual; para contrastar formalmente la normalidad se usan tests como Shapiro-Wilk o Kolmogorov-Smirnov.
  \end{exampleblock}
\end{frame}

%=============================================================
\section{Pruebas de normalidad}
%=============================================================

\begin{frame}{Test de Shapiro-Wilk}
  \begin{block}{Hipótesis}
    $H_0$: los datos provienen de una distribución Normal. \quad $H_1$: no Normal.
  \end{block}
  \bigskip
  \begin{itemize}
    \item Estadístico: $W = \dfrac{\left(\sum_{i=1}^n a_i x_{(i)}\right)^2}{\sum_{i=1}^n (x_i - \bar{x})^2}$, donde $a_i$ son coeficientes tabulados.
    \item $0 < W \le 1$. Cuanto más cerca de 1, más compatible con normalidad.
    \item \textbf{Recomendado} para $n \le 50$ (alta potencia).
  \end{itemize}
  \bigskip
  \begin{alertblock}{Regla de decisión}
    Si el $p$-valor $< \alpha$ (ej.\ 0.05) se rechaza $H_0$ (no hay normalidad).
  \end{alertblock}
\end{frame}

\begin{frame}{Test de Kolmogorov-Smirnov}
  \begin{block}{Estadístico}
    \[
      D = \sup_x |F_n(x) - F_0(x)|
    \]
    donde $F_n$ es la distribución empírica y $F_0$ la teórica.
  \end{block}
  \bigskip
  \begin{itemize}
    \item Versión de Lilliefors: cuando los parámetros de $F_0$ se estiman de los datos.
    \item Aplicable a cualquier distribución continua, no solo la Normal.
    \item Menor potencia que Shapiro-Wilk para detectar no-normalidad con muestras pequeñas.
  \end{itemize}
  \bigskip
  \begin{exampleblock}{Recomendación práctica}
    Usar Shapiro-Wilk para $n$ pequeño; Kolmogorov-Smirnov (Lilliefors) para $n$ grande.
  \end{exampleblock}
\end{frame}

%=============================================================
\section{Ajuste de distribuciones}
%=============================================================

\begin{frame}{Proceso de ajuste de una distribución teórica}
  \begin{enumerate}
    \item \textbf{Exploración visual}: histograma, boxplot, Q-Q plot.
    \item \textbf{Elección de candidatos}: según la forma y dominio de la variable.
    \item \textbf{Estimación de parámetros}: método de máxima verosimilitud (MV) o momentos.
    \item \textbf{Bondad de ajuste}: test $\chi^2$, Kolmogorov-Smirnov, AIC/BIC.
    \item \textbf{Selección final}: criterio estadístico + interpretabilidad.
  \end{enumerate}
  \bigskip
  \begin{block}{Distribuciones comunes en ciencias de la salud}
    Normal, Log-Normal, Exponencial, Gamma, Beta, Weibull, Binomial, Poisson.
  \end{block}
\end{frame}

%=============================================================
\section{Resumen}
%=============================================================

\begin{frame}{Resumen de la clase}
  \begin{block}{Herramientas vistas}
    \begin{itemize}
      \item \textbf{Histograma}: forma de la distribución, clases, densidad.
      \item \textbf{Boxplot}: $Q_1$, $Q_2$, $Q_3$, RIC, outliers.
      \item \textbf{Q-Q plot}: comparación con distribución teórica.
      \item \textbf{Tests}: Shapiro-Wilk, Kolmogorov-Smirnov.
    \end{itemize}
  \end{block}
  \bigskip
  \begin{block}{Próxima clase}
    Clase 10: Detección de outliers — criterios formales e informales, impacto en el análisis y estrategias de tratamiento.
  \end{block}
\end{frame}

\end{document}
